\section{The Higgs hypothesis}
\label{sec:higgs}

\subsection{The hypothesis and the conditions}

Hypothesis H2 says that the superfluid state, the condensate of the fermion pair along the hidden
axis, \emph{is the Higgs field} \cite{Higgs1964}. In the Standard Model the Higgs field is a doublet
of the weak-isospin group with a vacuum value $\vEW=246$~GeV, which breaks the electroweak symmetry,
gives the W and Z bosons and the quarks and charged leptons their masses, and leaves one scalar
particle of 125~GeV. Its four real components form a vector of a group $SO(4)$, which the vacuum
value breaks to $SO(3)$ \cite{Sikivie1980}. From these facts follow seven conditions that
a condensate must meet in order to \emph{be} the Higgs field (Table~\ref{tab:conditions}).
\SIG{With static additional times the hypothesis has no geometric basis on the $4+4$ spacetime,
because the condensate is a singlet under the only non-Abelian symmetry that the geometry
provides.}{With static extra coordinates the hypothesis is unproven on the $7+1$ spacetime, with one
structural point in its favor: the order parameter has the representation and the breaking pattern
of the Higgs field under a symmetry of the spacetime, the rotations of the four extra axes.} This
section examines the conditions on the deflating spacetime.

Two statements must be kept apart. H2 is the hypothesis that the \emph{pair condensate} is the Higgs
field; the electroweak scale is then the scale of the condensate. HV, used in one row of
Table~\ref{tab:bounds}, is the different assumption that the vacuum value of the
\emph{Vorton-Scalar} sets the electroweak scale. HV is not examined here.

\begin{table*}
\caption{\label{tab:conditions}The conditions that a condensate must meet in order to be the Higgs
field, and their status on the deflating \SIG{$4+4$}{$7+1$} spacetime. The last row is a condition
that only a deflating spacetime imposes.\SIG{ A dash marks a condition that is not examined for this
spacetime, because C3 is not met there by any symmetry of the spacetime.}{}}
\begin{ruledtabular}
\begin{tabular}{p{0.32\textwidth}p{0.62\textwidth}}
condition & status \\
\hline
C1: a scalar under the Lorentz group of ordinary spacetime & met in the sector of fields that do not depend on the extra coordinates \\
C2: non-zero in empty space & open; reached by the deflation only outside the range of validity of the model \\
C3: a doublet of the weak-isospin group with a neutral component & \SIG{only under a hidden symmetry; a singlet under the rotations of the additional times}{only under a hidden symmetry (three different rates), or with $SO(3)$ alone as a symmetry of the spacetime (equal rates)} \\
C4: breaks $SO(4)$ to $SO(3)$ & \SIG{met in the sector only}{met in the sector; the spacetime itself has at most the $SO(3)$} \\
C5: scale 246~GeV, with a scalar particle of 125~GeV with the observed couplings & \SIG{---}{unfavorable (estimate of Sec.~\ref{sec:c567})} \\
C6: its couplings give the quarks and leptons their masses & \SIG{---}{not addressed} \\
C7: no further light or charged particles that would have been seen & \SIG{---}{additional particles are predicted} \\
new: the electroweak scale does not drift & requires $\gsum<1.16\times10^{-6}\,H_0/\nEW$ \\
\end{tabular}
\end{ruledtabular}
\end{table*}

\subsection{What the deflation leaves unchanged}

The pair operators, their transformation properties and the internal symmetry are properties of the
Dirac matrices of the frame, which the deflation does not change, and of the equations of the fields
that do not depend on the extra coordinates, which for the weighted fields are the flat equations.
By Sec.~\ref{sec:hidden} the four vector pair operators form the $(2,2)$ of
$SU(2)_A\times SU(2)_B$, the representation of a Higgs field, and a condensate along one axis breaks
$SO(4)$ to $SO(3)$. Hence, \emph{within the sector}, C1 holds as on flat spacetime: the pair along
the hidden axis is invariant under the rotations of $x1,x2,x3$, which remain symmetries of the
spacetime, and under the boosts of ordinary spacetime as they act on the fields of the sector. (The
boosts are no longer symmetries of the deflating spacetime as a whole; in the sector this shows only
in the dependence of the couplings on the time.) C3 and C4 hold in the sector as well, in the sense
that the pattern is right under its internal symmetry.

\subsection{What the deflation changes}

%%<1671>
On the flat $7+1$ spacetime the six generators of $SO(4)$ are rotations of the extra dimensions. This
is the structural point in favor of the hypothesis: a symmetry of extra dimensions is the
kind of symmetry from which a gauge interaction can arise, as in the theories of Kaluza and Klein
\cite{Witten1981}. By Sec.~\ref{sec:hidden} the deflating spacetime keeps none, one or three of the
six, and never one of the three that mix the hidden axis with a deflating coordinate.

For equal rates the part of $SO(4)$ that remains a symmetry of the \emph{spacetime} is $SO(3)$, the
rotations of the three deflating coordinates: precisely the subgroup that a condensate along $x0$
would leave unbroken.
The deflation thus breaks $SO(4)$ to $SO(3)$ \emph{explicitly}, at the level of the geometry, in
the same pattern in which the hypothesis needs the condensate to break it \emph{spontaneously}.
Three consequences follow. First, the three broken generators, which under H2 would correspond to
the massive weak bosons, are not symmetries of the spacetime at all; a gauge interaction for them
cannot arise from the geometry of the deflating spacetime. Second, in the sector the full $SO(4)$
nevertheless acts as an exact symmetry, but as a hidden one. This is the situation of the model on
the spacetime with four space-like and four time-like dimensions, where the six generators are never
rotations of the spacetime and the hypothesis has no geometric basis. \emph{The deflation puts the
$7+1$ model in the position of the $4+4$ model.} Third, the potential of the order parameter does not single out the hidden axis
(Sec.~\ref{sec:order}): within the sector the direction of the condensate is as free as on flat
spacetime. Outside the sector the axes differ, because a variation of the order parameter along a
deflating coordinate costs a gradient energy that grows as $1/b_t^{2}$; the hidden symmetry is exact
only for fields that are strictly independent of the extra coordinates.
%%</1671>
%%<1644>
On the flat $4+4$ spacetime the only non-Abelian internal symmetry that the geometry provides is
$SO(3)_T$, the rotations of the three additional times, and under it the condensate along the hidden
axis is a singlet, not a doublet: the pair operator along $x0$ commutes with the rotations of
$x5,x6,x7$. By Sec.~\ref{sec:hidden} the deflating spacetime
keeps $SO(3)_T$ only if the three rates are equal, and reduces it to one rotation, or to none,
otherwise. The conclusion is the same where the rates are equal and holds a
fortiori where they are not: \emph{the hypothesis has no geometric basis on the $4+4$ spacetime.}
The pattern of the Higgs vacuum is present in the sector, under a hidden symmetry, and nowhere else.
(On a spacetime with seven space-like dimensions and one time the six generators are rotations of
the extra dimensions when these are static; a deflation of three of them removes that geometric
basis as well.)
%%</1644>

\subsection{Condition C2: a condensate in empty space}

With static extra coordinates C2 is open: a pair condensate forms in the vacuum only if $\aVS$
exceeds a critical value $\alpha_c$ of about 2 or more, depending on the cutoff (Sec.~\ref{sec:cjt}). On a deflating spacetime
$\aVS$ grows as $e^{\td}$ and reaches any such value at a definite time. Section~\ref{sec:cjt} showed
what else has happened by then: the exchange of one boson is no longer a controlled approximation,
and the field at which the one-loop potential of the Vorton-Scalar falls below its vacuum value has
moved to within a factor of two of the vacuum value. \emph{C2 can be reached by waiting, but not
within the range of validity of the model.} Its status remains open.

%%<1671>
\subsection{Conditions C5, C6 and C7: scale, masses and spectrum}
\label{sec:c567}

These conditions concern dynamics that the model does not determine, and the deflation does not
change what can be said about them.

\emph{C5.} In models in which a condensate of fermions replaces the elementary Higgs field
\cite{NJL1961,BardeenHillLindner1990} the scalar bound state that accompanies the condensate has a
mass of about twice the mass that the condensate gives to its constituents. With 125~GeV for the
scalar the constituents would weigh about 60~GeV. By C3 they carry electroweak charges, and
searches at the electron-positron collider LEP exclude new charged fermions below about 100~GeV
\cite{PDG2024}. The estimate speaks against the hypothesis; it is not a proof, because the dynamics
at strong coupling is not under control.

\emph{C6.} The model contains no coupling of the pair field to ordinary fermions. A pair of fermion
number 2 cannot couple to a fermion-antifermion bilinear of ordinary matter without violating the
fermion number or the conservation of the ordinary quantum numbers.

\emph{C7.} Besides the Higgs particle the model predicts the partners that complete the complex
$(2,2)$, a Goldstone boson of the fermion number, the fermions themselves with electroweak charges,
and the two modes of the Vorton-Scalar. None of these is observed.
%%</1671>

\subsection{The self-interaction and the hidden symmetry}
\label{sec:higgsself}

The internal symmetry of the sector, on which C3 and C4 rest, was established in
Sec.~\ref{sec:hidden} for the model without self-interaction. By Sec.~\ref{sec:selfsym} a scalar
and a pseudoscalar contact coupling keep the whole of it, and so does the exchange of the
Vorton-Scalar, which is a self-interaction of these two kinds with a finite range
[Eq.~(\ref{eq:gq_exchange})]. A vector or an axial-vector contact coupling keeps the six generators
of $SU(2)_A\times SU(2)_B$ and the phase symmetry but breaks $U(1)_E$. The representation and the
breaking pattern that C3 and C4 require are therefore present in the sector for any couplings of the
self-interaction; what a vector or axial-vector coupling removes is the additional symmetry
$U(1)_E$.

\subsection{A new condition: the electroweak scale must not drift}

Suppose that H2 holds. The electroweak scale $\vEW$ is then the scale of the pair condensate in the
vacuum, as in models in which a condensate of fermions replaces the elementary Higgs field
\cite{BardeenHillLindner1990}. That scale is determined by the parameters of the model: the masses
$m$ and $m_s$, the cutoff, and the coupling $\aVS$. On the deflating spacetime the masses are constant
and $\aVS\propto e^{\td}$. Therefore
\begin{equation}
\vEW\ \propto\ e^{\,\nEW\,\td},\qquad\nEW=\frac{d\ln\vEW}{d\ln\aVS}.
\label{eq:nEW}
\end{equation}
The exponent cannot vanish: below the critical coupling there is no condensate, so that the scale
depends on the coupling, and most strongly just above the critical coupling.

Under H2 the masses of the electron and of the quarks are proportional to $\vEW$ and change at the
rate $\nEW\,\gsum$. The mass of the proton is due mostly to the strong interaction and does not
follow. The ratio $m_p/m_e$ of the proton mass to the electron mass therefore changes at the rate
$\nEW\,\gsum$. Comparisons of atomic clocks find this ratio constant to $8.0\times10^{-17}$ per year
at two standard deviations \cite{Lange2021}. Hence
\begin{equation}
\gsum<\frac{8.0\times10^{-17}\ \text{per year}}{\nEW}=\frac{1.16\times10^{-6}\,H_0}{\nEW}.
\label{eq:clockbound}
\end{equation}
The bound needs no assumption about gravity. For $\nEW=1$ it is 2787 times stronger than the bound
that lunar laser ranging gives if gravity is eight-dimensional.

\emph{The size of the exponent.} For the model itself, with massive fermions and a massive radial
mode, the scale of the vacuum condensate has not been computed as a function of the coupling. It is
known in closed form for the scale-invariant case of the ladder equation: massless
fermions and massless bosons between an infrared scale and a cutoff \cite{Miransky1985}. There the
scale of the condensate obeys
\begin{equation}
\varrho_c\,\ln\frac{\text{cutoff}}{\text{scale}}=\arctan\frac{1}{2\varrho_c},\qquad
\varrho_c^{2}=\frac14\Big(\frac{\aVS}{\alpha_c}-1\Big),
\end{equation}
where $\varrho_c$ is only an abbreviation for the square root that the second equation defines, from
which
\begin{equation}
\nEW=\frac{(1+4\varrho_c^{2})\,\arctan(1/2\varrho_c)}{8\,\varrho_c^{3}}+\frac{1}{4\varrho_c^{2}}.
\label{eq:nEWladder}
\end{equation}
The program verifies Eq.~(\ref{eq:nEWladder}) against the numerical derivative to $3\times10^{-8}$
and tabulates it with the bound (\ref{eq:clockbound}) (Table~\ref{tab:higgs}). The exponent diverges
at the critical coupling, equals $1+\pi/2=2.57$ at twice the critical coupling, and exceeds 1 for
$\aVS<3.27\,\alpha_c$. It falls below 1 only where the scale of the condensate approaches the cutoff,
that is, where the four-dimensional description ends. Within the range in which a condensate of the
kind that H2 needs can be described at all, $\nEW$ is of order one or larger.

\begin{table}
\caption{\label{tab:higgs}The exponent $\nEW$ with which the scale of the vacuum condensate follows
the coupling in the scale-invariant ladder equation, and the bound (\ref{eq:clockbound}) on the
present rate of deflation that follows from it under hypothesis H2.}
\begin{ruledtabular}
\begin{tabular}{ccccc}
$\aVS/\alpha_c$ & $\ln$(cutoff/scale) & $\nEW$ & $\gsum$ (per year) & $\gsum/H_0$ \\
\hline
{{ROWS:tbl_higgs}}
\end{tabular}
\end{ruledtabular}
\end{table}

\emph{The conclusion.} Under H2 the deflation clock can have advanced by at most
$\gsum\,t_0<1.1\times10^{-6}/\nEW$ during the age of the universe at a constant rate. Every effect of
the deflation described in Secs.~\ref{sec:gap} to \ref{sec:cosmo} is then unobservably small.
Conversely, a deflation at any rate at which those effects matter excludes H2. \emph{The hypothesis
H2 and a deflation of any consequence today exclude each other.} This conclusion does not depend on
the unresolved conditions C2 and C5 to C7. It depends on two properties only: that the electroweak
scale is set by the dynamics of the model, which is the content of H2, and that the coupling of the
model follows the volume of the deflating coordinates, which is exact (Sec.~\ref{sec:clock}).

\subsection{Verdict on hypothesis H2}

On the deflating spacetime the hypothesis that the singlet superfluid is the Higgs field is
\SIG{unproven and without geometric basis, as on the flat $4+4$ spacetime,}{unproven, with less in
its favor than on the flat $7+1$ spacetime,} and it is incompatible with any deflation that is fast
enough to matter today. Whether a mechanism exists that would make the electroweak scale insensitive
to the coupling of the pair is not known; none is present in the model.

